\section{Problem and Design Goals}
 
The provenance gap decomposes into three problems. \emph{Provenance}: for any block of
randomness, which sources contributed, how much entropy each was credited, and when.
\emph{Integrity}: whether a source is still delivering the entropy it is credited with,
given that a failed source returns bytes of the correct type and raises no error.
\emph{Replay}: whether a third party can reconstruct the randomness a computation
consumed on different hardware and library versions, without trusting the original
machine.
 
Four failure modes with documented precedent motivate the design \citep{turan2018}:
silent source death, where a generator returns constants after a fault; over-extraction,
where a conditioner emits more bits than its input entropy justifies; over-trusted
assertion, where a source is credited at a rate it does not deliver; and implementation
drift, where a fixed seed yields different streams across library versions. An active
adversary who observes or influences the source at run time is out of scope; that is the
threat model of cryptographic generation, which the system does not address.
 
The design is organized around four goals, stated here as targets rather than as
guarantees, since their exact form will be settled by implementation.
 
\begin{itemize}
  \item \textbf{G1 (Budget accounting).} Conditioned output should be bounded by the
        min-entropy credited from sources, following the extraction bounds established
        in the theoretical literature, so that the system never silently degrades into a
        pseudo-random generator while presenting itself as physical.
  \item \textbf{G2 (Health monitoring).} Sources should be monitored continuously by
        standardized health tests, excluded from crediting when they fail, and auditable
        against their asserted rates by empirical estimation.
  \item \textbf{G3 (Replay).} Emitted randomness should be recorded with its provenance
        in a tamper-evident form, such that a computation can be re-run on the recorded
        randomness rather than on live sources.
  \item \textbf{G4 (Transparency).} Existing consumers of Python's and NumPy's random
        interfaces should operate unchanged, with the provenance layer confined to the
        seeding step and adding no statistical distortion.
\end{itemize}
 
Three exclusions are explicit: the system is not a cryptographic generator and must not
be used for keys or tokens; it does not stream physical entropy into every variate but
draws it to seed an expander; and the initial system is pure Python, with hardware
instructions and sensor sources as optional extras.
 